%======================================================================
\chapter{Arcs, cut points and simple closed curves}\label{ch:arcs}
%======================================================================

Everything in this chapter is general topology: no plane and no axioms of Moore's enter. We
develop the facts about arcs and simple closed curves that the rest of the volume uses constantly:
sub-arcs, first points, gluing, the characterization of arcs as irreducible continua and as
continua cut by each of their interior points, the characterization of simple closed curves as
continua cut by every pair of points, limits of sequences of arcs, and the replacement of a null
family of sub-arcs of an arc.

Throughout, $X$ is a Hausdorff space and $I=[0,1]$.

\section{Arcs and simple closed curves}

\begin{definition}\label{def:arc}
Let $A\ne B$ be points of $X$. A set $C\subseteq X$ is an \emph{arc from $A$ to $B$} if there is an
embedding $f\colon I\to X$ (a homeomorphism onto its image) with $f(I)=C$, $f(0)=A$ and $f(1)=B$.
$A$ and $B$ are the \emph{ends} of $C$; $\oarc C=C\setminus\set{A,B}$ is the \emph{open arc}. A
\emph{simple closed curve} is a set $J=C_1\cup C_2$ where $C_1,C_2$ are arcs from $A$ to $B$ with
$C_1\cap C_2=\set{A,B}$.
\end{definition}

\formal{\lean{HPC.Moore.Gen.TArc}, \lean{HPC.Moore.Gen.TSCC}.}

Since $I$ is compact and $X$ Hausdorff, a continuous injective $f\colon I\to X$ is automatically an
embedding (a continuous bijection from a compact space onto a Hausdorff space is a homeomorphism);
so to produce an arc it suffices to produce a continuous injection (\lean{HPC.Moore.Gen.TArc.of_inj}).
An arc from $A$ to $B$ is an arc from $B$ to $A$ (use $t\mapsto f(1-t)$), and it is compact,
closed and connected.

\begin{lemma}[Sub-arcs]\label{lem:subarc}
Let $f\colon I\to X$ be an embedding.
\begin{enumerate}[label=(\alph*)]
\item For $s<t$ in $I$, $f([s,t])$ is an arc from $f(s)$ to $f(t)$.
\item If $C$ is an arc from $A$ to $B$ and $P\ne Q$ are points of $C$, there is an arc from $P$ to
$Q$ contained in $C$.
\item The open arc $\oarc C$ of an arc $C$ is nonempty, connected and dense in $C$; so are
$C\setminus\set A$ and $C\setminus\set B$.
\end{enumerate}
\end{lemma}

\begin{proof}
(a) Compose $f$ with the affine homeomorphism $u\mapsto s+u(t-s)$ of $I$ onto $[s,t]$. (b) Write
$P=f(s)$, $Q=f(t)$ and apply (a) to $[\min(s,t),\max(s,t)]$, reversing if $s>t$. (c)
$\oarc C=f((0,1))$ is the image of a nonempty connected set, dense in $I$; similarly for
$f((0,1])$ and $f([0,1))$.
\end{proof}

\formal{\lean{HPC.Moore.Gen.TArc.image_Icc}, \lean{HPC.Moore.Gen.TArc.exists_sub},
\lean{HPC.Moore.Gen.TArc.nonempty_diff}, \lean{HPC.Moore.Gen.TArc.isPreconnected_diff},
\lean{HPC.Moore.Gen.TArc.subset_closure_diff}.}

The next lemma is used dozens of times in Chapters~\ref{ch:curves}--\ref{ch:quotient}.

\begin{lemma}\label{lem:conn-frontier}
Let $C,U\subseteq X$ with $C$ connected and $C\cap\fr U=\emptyset$. If $C$ meets $U$ then
$C\subseteq U$; if $C$ meets $X\setminus\cl U$ then $C\subseteq X\setminus\cl U$.
\end{lemma}

\begin{proof}
$C$ lies in $X\setminus\fr U=\interior U\cup(X\setminus\cl U)$, a union of two disjoint open sets,
so it lies in one of them.
\end{proof}

\formal{\lean{HPC.Moore.Gen.conn_sub_of_frontier}, \lean{HPC.Moore.Gen.conn_sub_ext_of_frontier}.}

\begin{lemma}[First points, gluing, joining]\label{lem:arc-ops}
\leavevmode
\begin{enumerate}[label=(\alph*)]
\item \emph{(First point.)} Let $C$ be an arc from $A$ to $B$ and $Z$ a closed set meeting $C$ with
$A\notin Z$. There are $P\in Z\cap C$ and an arc $AP\subseteq C$ with $AP\cap Z=\set P$.
\item \emph{(Gluing.)} If $D$ is an arc from $P$ to $Q$, $E$ an arc from $Q$ to $R$ and
$D\cap E=\set Q$, then $D\cup E$ is an arc from $P$ to $R$.
\item \emph{(Joining.)} If $\alpha$ is an arc from $P$ to $Q$, $\beta$ an arc from $R$ to $T$,
$\alpha\cap\beta\ne\emptyset$ and $P\ne R$, then $\alpha\cup\beta$ contains an arc from $P$ to
$R$.
\end{enumerate}
\end{lemma}

\begin{proof}
(a) Let $f$ parametrize $C$. The set $f^{-1}(Z)$ is closed, nonempty and misses $0$; let $t$ be its
least element and $P=f(t)$. Then $AP=f([0,t])$ works. (b) With $f$ for $D$ and $g$ for $E$, the map
$h(u)=f(2u)$ for $u\le\frac12$, $h(u)=g(2u-1)$ for $u\ge\frac12$ is continuous and injective (the
two halves only share $Q=h(\frac12)$). (c) If $P\in\beta$, take the sub-arc of $\beta$ from $P$ to
$R$. Otherwise apply (a) to $\alpha$ and the closed set $\beta$: we get $X_0\in\beta$ and an arc
$PX_0\subseteq\alpha$ meeting $\beta$ only in $X_0$. If $X_0=R$ we are done; otherwise glue $PX_0$
to the sub-arc $X_0R$ of $\beta$ by (b).
\end{proof}

\formal{\lean{HPC.Moore.Gen.TArc.first_point}, \lean{HPC.Moore.Gen.TArc.union},
\lean{HPC.Moore.Gen.TArc.join}.}

\section{Irreducible continua are arcs}

A set $C$ containing $A\ne B$ is \emph{irreducible between $A$ and $B$} if $C$ is connected and no
connected proper subset of $C$ contains both $A$ and $B$ (\vthree{Definition~2.1}). Every arc is
irreducible between its ends: a connected subset of $f(I)$ containing $f(0)$ and $f(1)$ is $f$ of a
connected subset of $I$ containing $0$ and $1$, which is all of $I$. The converse, in metrizable
spaces, is the main theorem of this section. In Moore's 1916 paper an \emph{arc} is \emph{defined}
as a bounded closed irreducible connected set (following Lennes), and the theorem identifies
Moore's arcs with the topological ones.

Fix a compact connected separable Hausdorff space $Y$ irreducible between $a\ne b$, i.e.\ every
connected $H\subseteq Y$ containing $a$ and $b$ is $Y$.

\begin{lemma}\label{lem:insert-conn}
Let $Z$ be a connected space, $p\in Z$, and $Z\setminus\set p=U\cup V$ with $U,V$ disjoint and open.
Then $U\cup\set p$ is connected.
\end{lemma}

\begin{proof}
$U\cup\set p=Z\setminus V$ is closed. If $U\cup\set p=F_1\cup F_2$ with $F_1,F_2$ disjoint, nonempty
and closed, say $p\in F_1$, then $F_2\subseteq U$ is closed in $Z$, and $F_2=U\setminus F_1$ is open
in $Z$. So $F_2$ is a nonempty proper clopen subset of $Z$, which is impossible.
\end{proof}

\formal{\lean{HPC.Moore.ac_insert_isPreconnected}.}

\begin{lemma}[Cuts]\label{lem:cut}
\leavevmode
\begin{enumerate}[label=(\alph*)]
\item $Y\setminus\set a$ and $Y\setminus\set b$ are connected.
\item For every $p\ne a,b$ there are disjoint connected open sets $U\ni a$, $V\ni b$ with
$Y\setminus\set p=U\cup V$.
\end{enumerate}
\end{lemma}

\begin{proof}
The basic observation is: \emph{if $Y\setminus\set p=U_2\cup W$ with $U_2,W$ disjoint and open,
$U_2\ne\emptyset$, then $W\cup\set p$ does not contain both $a$ and $b$}; for $W\cup\set p$ is
connected (Lemma~\ref{lem:insert-conn}) and is a proper subset of $Y$.

(a) If $Y\setminus\set b$ were not connected, write $Y\setminus\set b=U_2\cup W$ (disjoint, open,
nonempty), with $a\in W$. Then $W\cup\set b$ contains $a$ and $b$, contradicting the observation.

(b) $Y\setminus\set p$ is a proper subset of $Y$ containing $a$ and $b$, so it is not connected:
$Y\setminus\set p=U\cup V$ with $U,V$ disjoint, open and nonempty. By the observation (applied with
$(U_2,W)=(U,V)$ and with $(U_2,W)=(V,U)$) neither $U\cup\set p$ nor $V\cup\set p$ contains both $a$
and $b$; so $a$ and $b$ lie in different pieces, say $a\in U$, $b\in V$. If $U$ were not connected,
$U=U_1\cup U_2$ with $U_1,U_2$ disjoint, open and nonempty, $a\in U_1$; then $Y\setminus\set
p=U_2\cup(U_1\cup V)$ and $U_1\cup V\cup\set p$ contains $a$ and $b$, against the observation. So $U$
is connected, and likewise $V$.
\end{proof}

\formal{\lean{HPC.Moore.ac_compl_isPreconnected}, \lean{HPC.Moore.ac_cut}.}

For $x,y\in Y$ let $L(x,y)$ be the component of $Y\setminus\set y$ containing $x$ (empty if $x=y$).
Define $x\prec y$ if $x\in L(a,y)$, i.e.\ $x\ne y$ and $x$ lies on the $a$-side of $y$.

\begin{proposition}\label{prop:arc-order}
$\prec$ is a strict linear order on $Y$ with least element $a$ and greatest element $b$. The rays
$\set{x:x\prec y}=L(a,y)$ and $\set{x:y\prec x}=L(b,y)$ are open, and the order is dense: if $x\prec
y$ there is $z$ with $x\prec z\prec y$.
\end{proposition}

\begin{proof}
We use two facts. (i) By Lemma~\ref{lem:cut}, for $y\ne a,b$ the set $Y\setminus\set y$ is the
union of the disjoint connected open sets $L(a,y)$ and $L(b,y)$ (the sets $U,V$ of the lemma are
connected, so they are these components); $Y\setminus\set b=L(a,b)$ and $Y\setminus\set a=L(b,a)$.
In particular $b\notin L(a,y)$ for every $y$, and $L(a,a)=\emptyset$. (ii) A connected set
containing $x$ and missing $z$ lies in $L(x,z)$. By (i) and Lemma~\ref{lem:insert-conn}, the sets
$A_y=L(a,y)\cup\set y$ ($y\ne a$) and $B_y=L(b,y)\cup\set y$ ($y\ne b$) are connected.

\emph{Asymmetry.} Suppose $x\prec y$ and $y\prec x$. By (i), $x,y\notin\set{a,b}$. The connected set
$B_y$ contains $b$ and $y$ and misses $x$ (as $x\in L(a,y)$), so by (ii) $B_y\subseteq L(y,x)$,
which is $L(a,x)$ because $y\in L(a,x)$. Then $b\in L(a,x)$, against (i).

\emph{Transitivity.} Let $x\prec y\prec z$. Then $y\ne a$ (as $L(a,a)=\emptyset$), so $A_y$ is
connected; it contains $a$ and $x$, and misses $z$ ($z\ne y$, and $z\in L(a,y)$ would mean $z\prec
y$, against asymmetry). By (ii), $A_y\subseteq L(a,z)$, so $x\prec z$.

\emph{Totality.} Let $y\ne z$ with $y\not\prec z$. If $z=a$, then $a\in L(a,y)$, i.e.\ $z\prec y$.
The case $z=b$ is impossible, since $L(a,b)=Y\setminus\set b\ni y$. Otherwise $y\in L(b,z)$ by (i),
so $A_z$ is a connected set containing $a$ and $z$ and missing $y$, and $A_z\subseteq L(a,y)$ by
(ii): $z\prec y$.

\emph{Rays.} $x\in L(b,y)$ if and only if $y\prec x$: if $x\in L(b,y)$ then $x\ne y$ and $x\notin
L(a,y)$, so $y\prec x$ by totality; if $y\prec x$ then $y\ne b$, $x\ne y$ and $x\not\prec y$, so
$x\in L(b,y)$ by (i). Hence both rays are open. For $x\ne a$, $a\in L(a,x)$, i.e.\ $a\prec x$; for
$x\ne b$, $b\in L(b,x)$, i.e.\ $x\prec b$.

\emph{Density.} If $x\prec y$ and no $z$ lies between, then $Y$ is the union of
$\set{w:w\preceq x}=Y\setminus L(b,x)$ and $\set{w:y\preceq w}=Y\setminus L(a,y)$, two disjoint
nonempty closed sets, against connectedness.
\end{proof}

\formal{\lean{HPC.Moore.acL}, \lean{HPC.Moore.acL_asymm}, \lean{HPC.Moore.acL_trans},
\lean{HPC.Moore.acL_total}, \lean{HPC.Moore.acL_ray}, \lean{HPC.Moore.ac_exists_between}.}

\begin{theorem}[Irreducible continua are arcs]\label{thm:irr-arc}
A compact connected separable Hausdorff space $Y$ irreducible between $a\ne b$ is homeomorphic to
$I$ by a homeomorphism sending $0\mapsto a$, $1\mapsto b$.
\end{theorem}

\begin{proof}
Let $D\subseteq Y$ be countable and dense, and $D'=D\setminus\set{a,b}$. With the order of
Proposition~\ref{prop:arc-order}, $D'$ is a countable dense linear order without end points (between
two points of $Y$ there is a nonempty open interval, which meets $D$). It is nonempty, since
$Y\ne\set{a,b}$ (a two-point Hausdorff space is not connected). By Cantor's back-and-forth theorem
there is an order isomorphism $\varphi\colon D'\to\mathbb Q\cap(0,1)$. Define
$F(y)=\sup\set{\varphi(d):d\in D',\ d\prec y}$ (with $\sup\emptyset=0$). Then $F$ is
order preserving, strictly increasing (between $x\prec y$ there are two points of $D'$), $F(a)=0$
and $F(b)=1$. The preimage of an open ray $\set{t<s}$ or $\set{t>s}$ is a union of open rays of
$Y$, so $F$ is continuous. $F(Y)$ is a connected subset of $I$ containing $0$ and $1$, so
$F$ is onto. A continuous bijection from a compact space onto a Hausdorff space is a homeomorphism;
take its inverse.
\end{proof}

\formal{\lean{HPC.Moore.ac_exists_strictMono}, \lean{HPC.Moore.ac_homeomorph}.}

\begin{corollary}[Arc characterization]\label{cor:arc-char}
Let $X$ be metrizable, $A\ne B$ in $X$ and $C\subseteq X$. Then $C$ is an arc from $A$ to $B$ if and
only if $C$ is compact, connected, contains $A$ and $B$, and no connected proper subset of $C$
contains both $A$ and $B$.
\end{corollary}

\begin{proof}
One direction was noted above. Conversely, $C$ with the subspace topology is compact metrizable,
hence separable, and Theorem~\ref{thm:irr-arc} applies.
\end{proof}

\formal{\lean{HPC.Moore.isIrreducibleConnected_iff_arc}.}

\section{Cut-point arcs}

\begin{definition}\label{def:cutarc}
A set $H$ is a \emph{cut-point arc from $A$ to $B$} if $A\ne B$, $H$ is compact and connected,
$A,B\in H$, and $H\setminus\set x$ is not connected for every $x\in H\setminus\set{A,B}$.
\end{definition}

\formal{\lean{HPC.Moore.IsCutArc}.}

This is Moore's definition of an arc in his 1920 and 1925 papers (Chapter~\ref{ch:elements}).

\begin{proposition}\label{prop:cutarc}
Let $X$ be a $T_1$ space.
\begin{enumerate}[label=(\alph*)]
\item A cut-point arc from $A$ to $B$ is irreducible between $A$ and $B$.
\item If $X$ is metrizable, a cut-point arc from $A$ to $B$ is an arc from $A$ to $B$. Conversely,
the image of an embedding $f\colon I\to X$ is a cut-point arc from $f(0)$ to $f(1)$, and so is
$f([s,t])$ for $s<t$.
\item If $f\colon I\to X$ is an embedding and $L\subseteq f(I)$ is closed, connected and has two
points, then $L=f([s,t])$ for some $s<t$. If $f([s,t])$ is a cut-point arc from $a$ to $b$, then
$\set{a,b}=\set{f(s),f(t)}$.
\item \emph{(Nadler 6.7.)} If $K$ is a continuum and $H\subsetneq K$ is connected, some $x\in
K\setminus H$ has $K\setminus\set x$ connected.
\end{enumerate}
\end{proposition}

\begin{proof}
(d) is \vthree{Theorem~3.7}.
(a) Let $L\subsetneq H$ be connected with $A,B\in L$, and take $x\in H\setminus L$. By (d) applied to
the continuum $H$ and the connected set $L$ we may take $x$ with $H\setminus\set x$ connected; then
$x\ne A,B$ (as $A,B\in L$), against the definition.
(b) By (a) and Corollary~\ref{cor:arc-char}. For the converse: if $x=f(u)$ with $s<u<t$, then
$f([s,t])\setminus\set x=f([s,u))\cup f((u,t])$, two disjoint nonempty sets, each closed in the
union (they are compact images); so it is not connected.
(c) $f^{-1}(L)$ is a compact connected subset of $I$ with two points, i.e.\ an interval $[s,t]$,
$s<t$. In $f([s,t])$ removing $f(s)$ or $f(t)$ leaves a connected set, and removing any other point
does not; so the two points that can be removed without disconnecting are the ends.
\end{proof}

\formal{(a) \lean{HPC.Moore.IsCutArc.irreducible}; (b) \lean{HPC.Moore.IsCutArc.exists_embedding},
\lean{HPC.Moore.isCutArc_range}, \lean{HPC.Moore.isCutArc_image_Icc}; (c)
\lean{HPC.Moore.eq_image_Icc_of_subset_range}, \lean{HPC.Moore.ends_of_image_Icc}; (d)
\lean{HPC.Moore.exists_nonCut_diff}.}

\section{Simple closed curves}

\begin{theorem}[Moore 1920, Theorem~4]\label{thm:circle-char}
Let $J$ be a compact connected subset of a Hausdorff space with at least two points, such that
$J\setminus\set{x,y}$ is not connected for any two distinct $x,y\in J$. Then for any distinct
$p,q\in J$ there are cut-point arcs $a,b$ from $p$ to $q$ with $a\cup b=J$ and $a\cap b=\set{p,q}$.
\end{theorem}

\begin{proof}
Work in the space $Y=J$, which is compact, connected and Hausdorff.

\emph{Step 1: no point of $Y$ is a cut point.} Suppose $Y\setminus\set x=U\cup V$ with $U,V$
disjoint, open and nonempty. By Lemma~\ref{lem:insert-conn}, $U\cup\set x=Y\setminus V$ is a
continuum with at least two points, so it has two non-cut points (\vthree{Theorem~3.6}), one of
them, $u$, different from $x$; so $u\in U$ and $(U\cup\set x)\setminus\set u$ is connected.
Likewise there is $v\in V$ with $(V\cup\set x)\setminus\set v$ connected. Then
$Y\setminus\set{u,v}$ is the union of these two connected sets, which share $x$; so it is
connected, against the hypothesis.

\emph{Step 2: the halves.} $Y\setminus\set{p,q}$ is nonempty (otherwise $\set q=Y\setminus\set p$
would be a proper nonempty clopen set) and not connected, so $Y\setminus\set{p,q}=U\cup V$ with
$U,V$ disjoint, open and nonempty. Put $a=Y\setminus V=U\cup\set{p,q}$ and $b=Y\setminus
U=V\cup\set{p,q}$. Then $a,b$ are closed, $a\cup b=Y$ and $a\cap b=\set{p,q}$. We show that $a$ is
a cut-point arc from $p$ to $q$; the argument for $b$ is the same.

\emph{$a$ is connected.} Suppose $a=N\cup K$ with $N,K$ disjoint, closed and nonempty. If $p,q\in
N$, then $K\subseteq U$, so $K=U\setminus N$ is open in $Y$, and $K$ is a nonempty proper clopen
subset of $Y$: impossible. If $p\in K$ and $q\in N$, then one of the two pieces has a point other
than its marked point: if $K=\set p$ then $N\supseteq U\ne\emptyset$ has points other than $q$.
Suppose $K$ has a point $k\ne p$ (the other case is symmetric, exchanging $p$ and $q$). Then
\[
Y\setminus\set p=(K\setminus\set p)\cup\big((N\cup b)\setminus\set p\big),
\]
both sets are closed in $Y\setminus\set p$, they are disjoint (since $K\cap b\subseteq a\cap
b=\set{p,q}$ and $q\in N$), and nonempty ($k$, resp.\ $q$). This contradicts Step~1.

\emph{Each $x\in U$ cuts $a$.} Suppose $a\setminus\set x$ were connected. Take $y\in V$; we show
$S=Y\setminus\set{x,y}$ is connected, contradicting the hypothesis. Let $S\subseteq u\cup v$ with
$u,v$ open and $S\cap u\cap v=\emptyset$. The connected set $a\setminus\set x\subseteq S$ lies in
one of them, say $u$. Put $Q=S\cap v$, an open set. Every point of $S$ outside $V$ lies in
$a\setminus\set x\subseteq u$, so $Q\subseteq V$. If $z\ne y$ is in $\cl Q$, then $z\notin u$
(as $Q$ misses the open set $u$), $z\in\cl V\subseteq b$, so $z\ne x$ (as $x\in U$ misses $b$); thus
$z\in S\setminus u\subseteq v$ and $z\in Q$. So $Q$ is closed and open in $Y\setminus\set y$, which
is connected by Step~1. As $p\in Y\setminus\set y$ and $p\notin Q$, $Q=\emptyset$. Hence
$S\subseteq u$, and $S$ is connected.

The points of $a\setminus\set{p,q}$ are the points of $U$; so $a$ is a cut-point arc.
\end{proof}

\formal{\lean{HPC.Moore.circle_noCut}, \lean{HPC.Moore.circle_half}, \lean{HPC.Moore.circle_char}.}

\begin{proposition}\label{prop:scc-cut}
Let $J=X_1\cup X_2$ be a simple closed curve ($X_1,X_2$ arcs from $p$ to $q$, $X_1\cap
X_2=\set{p,q}$). Then $J\setminus\set{x,y}$ is not connected for any distinct $x,y\in J$, and $J$
is homeomorphic to the circle.
\end{proposition}

\begin{proof}
Let $f$, $g$ parametrize $X_1$, $X_2$ from $p$ to $q$. If $x=f(s)$, $y=f(t)$ with $s<t$ lie on one
arc, the closed sets $f([s,t])$ and $f([0,s]\cup[t,1])\cup X_2$ cover $J$ and meet only in $\set{x,y}$,
and each contains a point of $J\setminus\set{x,y}$. If $x=f(s)\in X_1\setminus X_2$ and $y=g(t)\in
X_2\setminus X_1$, use $f([0,s])\cup g([0,t])$ and $f([s,1])\cup g([t,1])$. In both cases
$J\setminus\set{x,y}$ is the union of two disjoint nonempty relatively closed sets.

For the second claim, the map $h\colon[0,1]\to J$ equal to $f(2u)$ for $u\le\frac12$ and to
$g(2-2u)$ for $u\ge\frac12$ is continuous, onto, and injective except that $h(0)=h(1)=p$; it
induces a continuous bijection from the circle $[0,1]/\set{0\sim1}$ onto $J$, which is a
homeomorphism by compactness.
\end{proof}

\formal{\lean{HPC.Moore.Gen.tscc_cut}; \lean{HPC.Jones.glue_circle},
\lean{HPC.Jones.tscc_homeomorph_circle}.}

So in a metrizable space the following are equivalent for a set $J$ with at least two points:
$J$ is a simple closed curve; $J$ is homeomorphic to the circle; $J$ is a continuum cut by every
pair of its points (Theorem~\ref{thm:circle-char}, Propositions~\ref{prop:cutarc}(b)
and~\ref{prop:scc-cut}).

\section{Limits of arcs and replacement of sub-arcs}

Two constructions with infinitely many arcs are needed in Chapter~\ref{ch:curves}.

\begin{theorem}[Limit arcs]\label{thm:limit-arc}
Let $X$ be metrizable, $B\in X$, and let $\varepsilon_n$ ($n\ge0$) be arcs from $y_n$ to $y_{n+1}$
such that $\varepsilon_n\cap\varepsilon_{n+1}=\set{y_{n+1}}$, $\varepsilon_n\cap\varepsilon_m=\emptyset$
for $m\ge n+2$, $B\notin\varepsilon_n$, and the $\varepsilon_n$ converge to $B$ (every
neighbourhood of $B$ contains $\varepsilon_n$ for all large $n$). Then
$\tau=\set B\cup\bigcup_n\varepsilon_n$ is an arc from $y_0$ to $B$.
\end{theorem}

\begin{proof}
$\tau$ is compact: an open cover has a member containing $B$, which contains all but finitely many
$\varepsilon_n$, and the remaining ones are compact. The same argument shows that
$\set B\cup\bigcup_{m\ge k}\varepsilon_m$ is compact for every $k$. The sets
$\Gamma_n=\varepsilon_0\cup\dots\cup\varepsilon_n$ are arcs from $y_0$ to $y_{n+1}$ (gluing,
Lemma~\ref{lem:arc-ops}(b)); they increase, so their union is connected, and $\tau$, which lies in its
closure, is connected.

$\tau$ is irreducible between $y_0$ and $B$. Let $z\in\tau\setminus\set{y_0,B}$ and let $n$ be least
with $z\in\varepsilon_n$. Let $g$ parametrize the arc $\Gamma_{n+1}$ from $y_0$ to $y_{n+2}$; then
$z=g(c)$ with $0<c<1$ ($z\ne y_{n+2}$ since $\varepsilon_{n+2}\cap\varepsilon_n=\emptyset$). The sets
$F_1=g([0,c])$ and $F_2=g([c,1])\cup\set B\cup\bigcup_{m\ge n+2}\varepsilon_m$ are compact, cover
$\tau$, and $F_1\cap F_2=\set z$ ($\varepsilon_m$, $m\ge n+2$, meets $\Gamma_{n+1}$ at most in
$y_{n+2}=g(1)$). So $\tau\setminus\set z$ is the union of the disjoint relatively closed sets
$F_1\setminus\set z\ni y_0$ and $F_2\setminus\set z\ni B$, and every connected subset of $\tau$
containing $y_0$ and $B$ contains $z$. Hence such a set is all of $\tau$, and
Corollary~\ref{cor:arc-char} applies.
\end{proof}

\formal{\lean{HPC.Moore.Gen.tarc_of_seq}; with trimming of arcs that meet more than consecutively,
\lean{HPC.Moore.Gen.exists_tarc_of_seq}.}

In applications the arcs $\beta_n$ come with $\beta_n\cap\beta_m=\emptyset$ for $m\ge n+2$ only;
cutting each $\beta_n$ at its first point on $\beta_{n+1}$ (Lemma~\ref{lem:arc-ops}(a)) produces arcs
$\varepsilon_n\subseteq\beta_n$ satisfying the hypotheses, and so $\set B\cup\bigcup\beta_n$ contains
an arc from $x_0$ to $B$.

\begin{theorem}[Replacement of a null family of sub-arcs]\label{thm:arc-replace}
Let $X$ be a metric space and $f\colon I\to X$ an embedding. Let $(u_i,v_i)$, $i\in\iota$, be
pairwise disjoint open subintervals of $I$, and $Z_i$ arcs from $f(u_i)$ to $f(v_i)$ whose open
arcs are pairwise disjoint and disjoint from $f(I)$. Suppose the family is \emph{null}: for every
$\varepsilon>0$ only finitely many $Z_i$ contain a point at distance $\ge\varepsilon$ from
$f(u_i)$. Then
\[
D=f\Big(I\setminus\bigcup_i(u_i,v_i)\Big)\cup\bigcup_i Z_i
\]
is an arc from $f(0)$ to $f(1)$.
\end{theorem}

\begin{proof}
Let $z_i\colon I\to Z_i$ parametrize $Z_i$ from $f(u_i)$ to $f(v_i)$, and define $h\colon I\to X$
by $h(t)=z_i\big((t-u_i)/(v_i-u_i)\big)$ for $t\in(u_i,v_i)$ and $h(t)=f(t)$ for $t$ in no interval.
Then $h(I)=D$, and $h$ is injective by the disjointness hypotheses. It remains to show that $h$ is
continuous; then $D$ is an arc. On each closed interval $[u_i,v_i]$ the map $h$ is continuous (the
two definitions agree at $u_i,v_i$). Let $s_k\to t$; we show $h(s_k)\to h(t)$, and may assume that
$t$ is in no open interval and that every $s_k$ lies in some interval $(u_{i_k},v_{i_k})$ (the
$s_k$ in no interval satisfy $h(s_k)=f(s_k)\to f(t)$). If some index $i$ occurs infinitely often,
then $t\in\set{u_i,v_i}$ and along those $k$, $h(s_k)\to h(t)$ by continuity on $[u_i,v_i]$. So we
may assume that every index occurs only finitely often. Since the intervals are disjoint, only
finitely many have length $\ge\delta$, for every $\delta>0$; so $v_{i_k}-u_{i_k}\to0$, hence
$u_{i_k}\to t$ and $f(u_{i_k})\to f(t)$. By nullity, $\sup_{x\in Z_{i_k}}d(x,f(u_{i_k}))\to0$. As
$h(s_k)\in Z_{i_k}$, $h(s_k)\to f(t)=h(t)$.
\end{proof}

\formal{\lean{HPC.Moore.Gen.arc_replace}.}

\section*{Exercises}
\addcontentsline{toc}{section}{Exercises}

\begin{exercise}
Show that in Lemma~\ref{lem:arc-ops}(c) the hypothesis $P\ne R$ is needed, and that the arc from $P$
to $R$ need not be unique.
\end{exercise}

\begin{exercise}
Show that the conclusion of Corollary~\ref{cor:arc-char} fails for the ``long line'' compactified
by two end points (a compact connected Hausdorff space irreducible between its end points which is
not metrizable). Which step of the proof of Theorem~\ref{thm:irr-arc} fails?
\end{exercise}

\begin{exercise}
Let $J$ be a simple closed curve and $p,q\in J$ distinct. Show that there are exactly two arcs in
$J$ from $p$ to $q$, and that their union is $J$.
\end{exercise}

\begin{exercise}
Show that the null hypothesis in Theorem~\ref{thm:arc-replace} cannot be omitted: replace the
intervals $(1/(n+1),1/n)$ of the segment $[0,1]\times\set0$ by semicircles of radius $1$ (suitably
tilted so that they are disjoint) and show that the result is not compact.
\end{exercise}
